\documentclass[12pt,a4paper]{article}

\usepackage[spanish]{babel}
\usepackage[utf8]{inputenc}
\usepackage[T1]{fontenc}
\usepackage{geometry}
\usepackage{xcolor}
\usepackage{hyperref}
\usepackage{graphicx}
\usepackage{float}
\usepackage{array}
\usepackage{longtable}
\usepackage{booktabs}
\usepackage{enumitem}
\usepackage{fancyhdr}
\usepackage{listings}
\usepackage{caption}
\usepackage{tikz}
\usetikzlibrary{arrows.meta,positioning,shapes.geometric,fit}

\geometry{margin=2.5cm}
\setlength{\parskip}{0.65em}
\setlength{\parindent}{0pt}

\definecolor{EduGreen}{RGB}{18,126,107}
\definecolor{EduLight}{RGB}{238,247,244}
\definecolor{EduGray}{RGB}{245,247,247}
\definecolor{EduBorder}{RGB}{126,155,148}

\hypersetup{
    colorlinks=true,
    linkcolor=EduGreen,
    urlcolor=EduGreen,
    citecolor=EduGreen,
    pdftitle={MTD EduPKIManager},
    pdfauthor={EduPKIManager}
}

\pagestyle{fancy}
\fancyhf{}
\lhead{Proyecto 9}
\rhead{MTD EduPKIManager}
\cfoot{\thepage}

\lstdefinestyle{code}{
    basicstyle=\ttfamily\small,
    breaklines=true,
    frame=single,
    backgroundcolor=\color{EduGray},
    columns=fullflexible
}

\newcommand{\captura}[3]{
    \begin{figure}[H]
        \centering
        \IfFileExists{img_latex/#1.png}{
            \includegraphics[width=0.92\textwidth,height=0.58\textheight,keepaspectratio]{img_latex/#1.png}
        }{
            \fbox{
                \begin{minipage}[c][4.6cm][c]{0.88\textwidth}
                    \centering
                    \textbf{#2}\\[0.35cm]
                    \small Falta la imagen: \texttt{img\_latex/#1.png}
                \end{minipage}
            }
        }
        \caption{#3}
    \end{figure}
}

\newcommand{\diagrama}[2]{
    \begin{figure}[H]
        \centering
        \IfFileExists{img_latex/#1.png}{
            \includegraphics[width=0.96\textwidth,height=0.62\textheight,keepaspectratio]{img_latex/#1.png}
        }{
            \fbox{
                \begin{minipage}[c][5cm][c]{0.88\textwidth}
                    \centering
                    \textbf{#2}\\[0.35cm]
                    \small Falta la imagen: \texttt{img\_latex/#1.png}
                \end{minipage}
            }
        }
        \caption{#2}
    \end{figure}
}

\begin{document}

\begin{titlepage}
    \centering
    \vspace*{1.7cm}
    {\Huge\bfseries Universidad La Salle - Arequipa\par}
    \vspace{1cm}
    {\Large Facultad de Ingenieria\par}
    \vspace{0.5cm}
    {\Large Escuela Profesional de Ingenieria de Software\par}

    \vspace{3.0cm}
    {\Huge\bfseries MTD\par}
    \vspace{0.6cm}
    {\huge\bfseries EduPKIManager\par}
    \vspace{0.4cm}
    {\Large Plataforma completa de Infraestructura de Clave Publica educativa\par}

    \vspace{1.2cm}
    {\Large\textbf{Proyecto 9:} PKI, Certificados X.509, Firmas Digitales y Criptografia Asimetrica\par}

    \vfill
    {\Large\textbf{Curso:} Ciberseguridad\par}
    \vspace{0.35cm}
    {\Large\textbf{Semestre:} X\par}
    \vspace{0.35cm}
    {\Large\textbf{Docente:} Marco Antonio Camacho Alatrista\par}
    \vspace{0.35cm}
    {\Large\textbf{Estudiante:} [Nombre del estudiante]\par}

    \vspace{1.1cm}
    {\Large Arequipa, Peru\par}
    {\Large 12 de julio de 2026\par}
\end{titlepage}

\tableofcontents
\newpage

\section{Objetivo}

Desarrollar una plataforma completa de Infraestructura de Clave Publica que implemente una Autoridad Certificadora propia, gestion del ciclo de vida de certificados digitales, firmas digitales de documentos y soporte TLS 1.3, orientada a entornos empresariales educativos.

El sistema busca demostrar de forma integrada los conceptos de PKI, certificados X.509, criptografia asimetrica, validacion de confianza, revocacion, OCSP, firma PDF, auditoria y despliegue reproducible con Docker Compose.

\section{Titulo del Software}

\textbf{EduPKIManager - Plataforma PKI educativa}

\section{Tecnologias Utilizadas}

\begin{longtable}{>{\bfseries}p{4cm}p{10cm}}
\toprule
Categoria & Tecnologia \\
\midrule
Lenguaje backend & Python \\
Framework backend & Django REST Framework \\
Criptografia & \texttt{cryptography}, OpenSSL mediante herramientas externas y \texttt{pyHanko} para PAdES \\
Frontend & React.js con TypeScript y Vite \\
Base de datos & PostgreSQL \\
Servidor web & Nginx para frontend HTTPS y proxy local \\
Despliegue & Docker Compose \\
Protocolos & X.509 v3, CRL, OCSP RFC 6960, TLS 1.3 \\
Algoritmos & RSA-2048, RSA-4096, ECDSA P-256, SHA-256 \\
Auditoria & Registro JSONL append-only con hash encadenado \\
\bottomrule
\end{longtable}

\section{Requisitos Funcionales}

\begin{longtable}{>{\bfseries}p{2.4cm}p{8.2cm}p{3.2cm}}
\toprule
Codigo & Requisito & Estado \\
\midrule
RF-01 & Implementacion de Autoridad Certificadora raiz con generacion de par de claves RSA-4096. & Cumplido \\
RF-02 & Emision de certificados X.509 v3 para usuarios, servidores y dispositivos con campos estandar completos. & Cumplido \\
RF-03 & Gestion del ciclo de vida de certificados: emision, renovacion, suspension y revocacion. & Cumplido \\
RF-04 & Publicacion y consulta de Lista de Revocacion de Certificados actualizada automaticamente. & Cumplido \\
RF-05 & Soporte OCSP para verificacion en linea de estado. & Cumplido \\
RF-06 & Modulo de firma digital de documentos PDF con certificado emitido por la CA propia. & Cumplido \\
RF-07 & Verificacion de firmas digitales con trazabilidad completa de cadena de confianza. & Cumplido \\
RF-08 & Implementacion de handshake TLS 1.3 usando certificados emitidos por la CA del sistema. & Cumplido \\
RF-09 & Panel de administracion web para gestion completa de certificados con roles administrador y usuario final. & Cumplido \\
RF-10 & Registro de auditoria inmutable de todas las operaciones criptograficas realizadas en la plataforma. & Cumplido \\
\bottomrule
\end{longtable}

\section{Diagramas del Sistema}

\subsection{Diagrama de Componentes}

La arquitectura se organiza en cuatro capas principales: presentacion, API, servicios criptograficos y persistencia/artefactos. El frontend consume la API REST y el backend delega las operaciones criptograficas en modulos especializados.

\diagrama{diagrama_componentes_edupkimanager}{Diagrama de Componentes de EduPKIManager}

\subsection{Diagrama de Modulos}

El backend fue separado por responsabilidad. Cada modulo concentra una parte del dominio PKI para mantener la implementacion clara y verificable.

\diagrama{diagrama_modulos_backend}{Diagrama de Modulos del Backend}

\subsection{Diagrama de Flujo: Emision de Certificado}

La emision parte de una solicitud autenticada. Segun el rol, el backend valida permisos, genera la clave privada del certificado final, construye un certificado X.509 v3 y lo firma con la Intermediate CA.

\diagrama{diagrama_flujo_emision_certificado}{Diagrama de Flujo de Emision de Certificado}

\subsection{Diagrama de Flujo: Verificacion de Firma PDF}

La verificacion de firma valida integridad del documento, validez criptografica de la firma, cadena Root/Intermediate, estado de revocacion y proposito de uso.

\diagrama{diagrama_flujo_verificacion_pdf}{Diagrama de Flujo de Verificacion de Firma PDF}

\subsection{Diagrama de Secuencia: Handshake TLS 1.3}

El despliegue Docker genera un certificado de servidor para \texttt{edupkimanager.com}. Nginx sirve el frontend por HTTPS restringido a TLS 1.3 usando dicho certificado.

\diagrama{diagrama_secuencia_tls13}{Diagrama de Secuencia de Handshake TLS 1.3}

\section{Implementacion}

\subsection{Estructura del Proyecto}

\begin{lstlisting}[style=code]
EduPKIManager/
  backend/
    api/          # Endpoints REST, autenticacion y permisos
    audit/        # Auditoria append-only con hash encadenado
    ca/           # Root CA, Intermediate CA y certificados X.509
    crl/          # Publicacion de CRL versionada
    ocsp/         # Responder OCSP JSON y DER RFC 6960
    pdf_sign/     # Firma PDF detached y PAdES
    tls/          # Certificado servidor y contexto TLS 1.3
    validation/   # Validacion de confianza X.509
    scripts/      # Demos, smoke tests y evidencias
    tests/        # Pruebas unitarias
  frontend/
    src/
      admin/      # Panel administrador
      user/       # Portal usuario final
      tools/      # Firma PDF para usuarios; OCSP, CRL, TLS y auditoria para admin
  scripts/
    setup_windows_tls.ps1
  docker-compose.yml
\end{lstlisting}

\subsection{Autoridad Certificadora}

El modulo de CA genera una cadena de confianza de tres niveles:

\begin{lstlisting}[style=code]
EduPKIManager Root CA
  -> EduPKIManager Intermediate CA
      -> Certificados finales de usuario, servidor y dispositivo
\end{lstlisting}

La Root CA utiliza RSA-4096 y actua como ancla de confianza. La Intermediate CA firma certificados finales, CRL y respuestas OCSP. Las claves privadas de CA se almacenan cifradas con una contrasena definida por la variable \texttt{CA\_KEY\_PASSWORD}.

\subsection{Emision X.509 v3}

Los certificados emitidos incluyen campos y extensiones estandar:

\begin{itemize}[leftmargin=*]
    \item Subject y Issuer.
    \item Numero de serie unico.
    \item Periodo de vigencia.
    \item Subject Alternative Name.
    \item Basic Constraints.
    \item Key Usage.
    \item Extended Key Usage.
    \item Subject Key Identifier y Authority Key Identifier.
    \item CRL Distribution Points.
    \item Authority Information Access con OCSP.
    \item Certificate Policies.
\end{itemize}

\subsection{Gestion de Ciclo de Vida}

El ciclo de vida se implementa mediante endpoints REST y estados persistidos en PostgreSQL.

\begin{longtable}{>{\bfseries}p{3.5cm}p{10.5cm}}
\toprule
Accion & Comportamiento \\
\midrule
Emision & Genera clave privada, certificado X.509 v3 y registro persistente. \\
Renovacion & Marca el certificado anterior como renovado y emite uno nuevo. \\
Suspension & Cambia el estado a \texttt{suspended} y publica la razon \texttt{certificate\_hold}. \\
Revocacion & Cambia el estado a \texttt{revoked} y agrega el certificado a la CRL. \\
\bottomrule
\end{longtable}

\subsection{CRL Versionada}

La CRL se publica en PEM y DER. El servicio \texttt{crl\_scheduler} refresca automaticamente los artefactos y mantiene un manifiesto con versiones historicas.

\begin{lstlisting}[style=code]
data/crl/crl-1.pem
data/crl/crl-1.der
data/crl/latest.crl.pem
data/crl/latest.crl.der
data/crl/manifest.json
\end{lstlisting}

\subsection{OCSP}

El sistema ofrece dos formas de consulta:

\begin{itemize}[leftmargin=*]
    \item Endpoint JSON para uso del frontend: \texttt{POST /api/ocsp/status/}.
    \item Endpoint DER RFC 6960 para clientes estandar: \texttt{POST /api/ocsp/} y \texttt{POST /ocsp/}.
\end{itemize}

Los estados soportados son \texttt{good}, \texttt{revoked} y \texttt{unknown}.

\subsection{Firma y Verificacion PDF}

La firma desprendida calcula el hash SHA-256 del PDF, firma el digest con la clave privada asociada al certificado y devuelve un sobre JSON con los datos necesarios para verificacion. La firma PAdES genera un PDF firmado embebido usando \texttt{pyHanko}.

Durante la verificacion se comprueba:

\begin{itemize}[leftmargin=*]
    \item Integridad del documento.
    \item Firma criptografica.
    \item Certificado emisor.
    \item Cadena Root/Intermediate.
    \item Proposito de firma documental.
    \item Estado local, CRL y OCSP.
\end{itemize}

\subsection{TLS 1.3}

El servicio \texttt{tls\_init} genera un bundle de servidor para \texttt{edupkimanager.com}. Nginx monta el certificado y restringe el protocolo a TLS 1.3.

\begin{lstlisting}[style=code]
ssl_protocols TLSv1.3;
ssl_certificate /etc/nginx/edupki/server_cert.pem;
ssl_certificate_key /etc/nginx/edupki/server_key.pem;
\end{lstlisting}

\subsection{Auditoria Inmutable}

Cada operacion criptografica relevante se registra en formato JSONL con una cadena de hashes.

\begin{lstlisting}[style=code]
entry_hash = SHA256(payload_canonico + previous_hash)
\end{lstlisting}

El endpoint de auditoria recalcula la cadena y reporta si alguna entrada fue alterada o eliminada.

\section{API Principal}

\begin{longtable}{>{\bfseries}p{4.2cm}p{9.8cm}}
\toprule
Endpoint & Descripcion \\
\midrule
\texttt{GET /api/health/} & Verifica que la API este levantada. \\
\texttt{GET /api/readiness/} & Valida base de datos, CA, CRL y auditoria. \\
\texttt{POST /api/auth/login/} & Autentica administrador o usuario final y devuelve token. \\
\texttt{GET /api/ca/root.pem} & Descarga la Root CA publica. \\
\texttt{POST /api/certificates/} & Emite certificado X.509 v3. \\
\texttt{GET /api/certificates/} & Lista certificados segun rol. \\
\texttt{POST /api/certificates/\{id\}/renew/} & Renueva certificado. \\
\texttt{POST /api/certificates/\{id\}/suspend/} & Suspende certificado. \\
\texttt{POST /api/certificates/\{id\}/revoke/} & Revoca certificado. \\
\texttt{GET /api/crl.pem} & Descarga CRL PEM actual. \\
\texttt{GET /api/crl.der} & Descarga CRL DER actual. \\
\texttt{GET /api/crl/manifest/} & Lista versiones CRL publicadas. \\
\texttt{POST /api/ocsp/status/} & Consulta estado OCSP en JSON. \\
\texttt{POST /api/ocsp/} & Responder OCSP estandar DER. \\
\texttt{POST /api/pdf/sign/} & Firma PDF con sobre JSON desprendido. \\
\texttt{POST /api/pdf/sign-embedded/} & Firma PDF PAdES embebido. \\
\texttt{POST /api/pdf/verify/} & Verifica firma JSON desprendida. \\
\texttt{POST /api/pdf/verify-embedded/} & Verifica PDF PAdES. \\
\texttt{POST /api/certificates/validate/} & Valida confianza X.509 por proposito. \\
\texttt{GET /api/tls/demo/} & Ejecuta prueba de handshake TLS 1.3. \\
\texttt{GET /api/audit/} & Lista auditoria y valida hash encadenado. \\
\bottomrule
\end{longtable}

\section{Interfaz de Usuario}

\subsection{Login}

La pantalla de login separa el acceso por credenciales demo. Despues de autenticar, el frontend conserva el token y habilita las vistas de acuerdo con el rol.

\captura{login_edupkimanager}{Login de EduPKIManager}{Interfaz de inicio de sesion}

\subsection{Panel Administrador}

El panel administrador ofrece gestion completa del ciclo de vida: emision, listado, renovacion, suspension y revocacion. Tambien permite descargar artefactos y operar funciones avanzadas de PKI.

\captura{panel_administrador}{Panel administrador}{Interfaz de gestion completa de certificados}

\subsection{Portal Usuario Final}

El portal usuario final reduce el alcance a certificados propios emitidos por el administrador. Desde esta vista se puede ver el estado, descargar certificados asignados y utilizarlos para operaciones de firma y verificacion PDF.

\captura{portal_usuario}{Portal de usuario}{Interfaz del usuario final}

\subsection{Operaciones PKI}

La vista de operaciones agrupa herramientas segun rol. El administrador accede a firma PDF, verificacion PDF, confianza X.509, OCSP, CA/CRL, TLS 1.3 y auditoria. El usuario final accede solo a firma y verificacion PDF.

\captura{operaciones_pki}{Operaciones PKI}{Interfaz de herramientas criptograficas}

\section{Despliegue}

\subsection{Configuracion}

El proyecto incluye \texttt{.env.example}. Para un despliegue serio se debe copiar a \texttt{.env} y cambiar los secretos.

\begin{lstlisting}[style=code]
cp .env.example .env
python backend/scripts/check_deployment_config.py --env-file .env --profile production
\end{lstlisting}

\subsection{Ejecucion con Docker Compose}

\begin{lstlisting}[style=code]
docker compose up -d --build
docker compose ps
\end{lstlisting}

Servicios esperados:

\begin{itemize}[leftmargin=*]
    \item \texttt{db}: PostgreSQL.
    \item \texttt{backend}: API Django REST.
    \item \texttt{frontend}: Nginx + React.
    \item \texttt{crl\_scheduler}: publicacion periodica de CRL.
    \item \texttt{tls\_init}: generacion de certificados TLS.
\end{itemize}

\subsection{Dominio Local}

Para usar el dominio local:

\begin{lstlisting}[style=code]
127.0.0.1 edupkimanager.com www.edupkimanager.com
\end{lstlisting}

En Windows se automatiza con:

\begin{lstlisting}[style=code]
.\scripts\setup_windows_tls.ps1
.\scripts\setup_windows_tls.ps1 -TrustRootCa
\end{lstlisting}

\section{Pruebas y Validacion}

El proyecto incluye pruebas unitarias, pruebas de humo, demo end-to-end, verificacion TLS y generacion de evidencias.

\begin{lstlisting}[style=code]
PYTHONPATH=backend python -m unittest discover backend/tests
PYTHONPATH=backend python backend/scripts/e2e_demo.py
PYTHONPATH=backend python backend/scripts/tls13_demo.py
python backend/scripts/api_smoke_test.py http://127.0.0.1:8000/api
PYTHONPATH=backend python backend/scripts/deployment_acceptance.py
PYTHONPATH=backend/scripts python backend/scripts/generate_evidence_bundle.py
\end{lstlisting}

\subsection{Criterios de Aceptacion}

\begin{itemize}[leftmargin=*]
    \item La API responde \texttt{health} y \texttt{readiness}.
    \item La Root CA y la Intermediate CA existen.
    \item Se puede emitir, renovar, suspender y revocar certificados.
    \item La CRL se publica en PEM/DER y mantiene manifiesto.
    \item OCSP responde estados correctos.
    \item Se firma y verifica PDF en formato JSON y PAdES.
    \item TLS 1.3 responde para \url{https://edupkimanager.com}.
    \item La auditoria mantiene cadena valida.
\end{itemize}

\section{Seguridad}

\subsection{Controles Implementados}

\begin{itemize}[leftmargin=*]
    \item Cifrado de claves privadas de CA en reposo.
    \item Separacion Root CA e Intermediate CA.
    \item Autenticacion con token firmado.
    \item Control de roles administrador y usuario.
    \item Propiedad de certificados para evitar acceso cruzado.
    \item Validacion de uso criptografico por proposito.
    \item Publicacion de CRL y soporte OCSP.
    \item Auditoria append-only con hash encadenado.
    \item TLS 1.3 en el acceso HTTPS local.
\end{itemize}

\subsection{Consideraciones para Produccion}

\begin{itemize}[leftmargin=*]
    \item Cambiar todas las credenciales demo.
    \item Rotar secretos de forma controlada.
    \item Proteger los volumenes de CA y auditoria.
    \item Usar backups cifrados.
    \item Limitar acceso administrativo.
    \item Registrar eventos en un sistema externo si se requiere mayor resistencia.
\end{itemize}

\section{Repositorio y Control de Versiones}

El proyecto fue inicializado con Git y publicado en:

\begin{center}
\url{https://github.com/abeluciano/EduPKIManager.git}
\end{center}

Los commits se organizaron con Conventional Commits en espanol, separando configuracion, backend, frontend, scripts, pruebas y documentacion.

\section{Conclusiones}

EduPKIManager implementa una plataforma PKI completa para fines educativos. El sistema cubre generacion de CA, emision X.509 v3, ciclo de vida de certificados, CRL automatica, OCSP, firma y verificacion PDF, TLS 1.3, panel web por roles y auditoria inmutable. La solucion queda preparada para demostraciones reproducibles mediante Docker Compose y scripts de validacion.

La arquitectura modular facilita ampliar el proyecto con integraciones empresariales, autenticacion institucional, HSM, politicas de certificados avanzadas o monitoreo externo.

\end{document}
